\section{Order addition at every multiplicative depth}
\label{rlc:sec:convolution}

For $r\ge2$, set $x_r=\mu-x_1-\cdots-x_{r-1}$ and define
\begin{equation}\label{rlc:eq:convolution}
 (f_{a,s_1}*\cdots*f_{a,s_r})(\mu)
 =\int_{\R^{r-1}}\prod_{j=1}^r f_{a,s_j}(x_j)
       \dd x_1\cdots\dd x_{r-1}.
\end{equation}
Equivalently $y_j=e^{x_j}$ and $y_1\cdots y_r=e^\mu$, with measure
$\prod_{j<r}\dd y_j/y_j$. For $r=1$ use ordinary evaluation.

\begin{theorem}[Entire-order addition]\label{rlc:thm:addition}
For $\Rea a>0$, real $\mu$, and arbitrary $s_1,\ldots,s_r\in\C$,
\eqref{rlc:eq:convolution} is absolutely convergent and depends on the
orders only through $W=s_1+\cdots+s_r$. Denote it by $\CC_{r,a}(W;\mu)$.
It is entire in $W$ and has the representation
\begin{equation}\label{rlc:eq:barnes}
 \CC_{r,a}(W;\mu)=\frac1{2\pi i}\int_{c-i\infty}^{c+i\infty}
 e^{p\mu}\left(\frac{\pi}{\sin\pi p}\right)^r
 (p+a-1)^{-W}\dd p,
\end{equation}
where $c$ satisfies \eqref{rlc:eq:cstrip}. This formula extends it
holomorphically to $|\Ima\mu|<r\pi$.
Every fixed mixed order derivative can be passed through
\eqref{rlc:eq:convolution}; a derivative of total degree $M$ gives
$\partial_W^M\CC_{r,a}(W;\mu)$.
\end{theorem}
\begin{proof}
Put $g_j(x)=e^{-cx}f_{a,s_j}(x)$. The factor $e^{c\mu}$ is constant
on the constraint hyperplane. The integral of absolute values in
\eqref{rlc:eq:convolution} is therefore at most
\[
 e^{c\mu}\|g_r\|_\infty\prod_{j<r}\|g_j\|_1.
\]
Lemma~\ref{rlc:lem:weighted} proves absolute convergence and locally
uniform bounds for all parameter derivatives. The Fourier transform of
$g_j$ is \eqref{rlc:eq:transform} at $p=c+it$. Its product is
\[
 \left(\frac{\pi}{\sin\pi p}\right)^r(p+a-1)^{-W}.
\]
This depends only on $W$. Fourier inversion is applicable because the
product decays like $e^{-r\pi|t|}$ times a fixed power of $1+|t|$ on
compact order sets. It gives \eqref{rlc:eq:barnes} and establishes
pointwise equality, not merely equality almost everywhere; the
convolutions are continuous. The same exponential decay proves the
stated $\mu$ strip and entire dependence on $W$. Norm holomorphy and
Cauchy's formula justify all differentiated integrals.
\end{proof}

For $\Rea W>0$, another useful form follows from the ordinary
one-sided fractional-integral semigroup:
\begin{equation}\label{rlc:eq:basefractional}
 \CC_{r,a}(W;\mu)=\frac1{\Gamma(W)}
 \int_0^\infty u^{W-1}e^{-(a-1)u}
           f_0^{*r}(\mu-u)\dd u.
\end{equation}
One can prove it first with all $\Rea s_j>0$ by Fubini and the
Dirichlet simplex integral; their sum of positive integration variables
has Gamma density of order $W$. Alternatively its transform is
\eqref{rlc:eq:barnes}. The kernel in
\eqref{rlc:eq:basefractional} is then continued through $W=0,-1,\ldots$
by the already established entire function, not by substituting into an
improper integral that no longer converges.

\subsection{A direct two-factor proof}
It is useful to verify the main normalization without Fourier inversion.
For $\Rea s,\Rea t>0$, insert the two Euler integrals and integrate the
reciprocal variable first. For $u,v>0$,
\[
 \int_0^\infty
 \frac{\dd y/y}{(1+ye^{-u})(1+y^{-1}e^{-v})}
 =\frac{u+v}{1-e^{-(u+v)}}.
\]
This follows by one partial fraction decomposition; the logarithms at
infinity cancel. With $w=u+v$, the remaining beta integral yields
\begin{align*}
 \int_0^\infty F_{a,s}(y)F_{a,t}(1/y)\frac{\dd y}{y}
 &=\frac1{\Gamma(W)}\int_0^\infty
       \frac{w^W e^{-(a-1)w}}{e^w-1}\dd w\\
 &=W\zeta(W+1,a).
\end{align*}
Absolute convergence and entire-order continuation proved above give
\eqref{rlc:eq:opening} for all $s,t$. In particular,
\begin{equation}\label{rlc:eq:polylogpair}
 \int_0^\infty\Li_s(-y)\Li_t(-1/y)\frac{\dd y}{y}
     =(s+t)\zeta(s+t+1),
\end{equation}
with the removable value at zero understood. For example, the integrals
for the order pairs $(1,1)$, $(2,1)$, and $(2,2)$ are respectively
$2\zeta(3)$, $3\zeta(4)$, and $4\zeta(5)$.
These examples illustrate the normalization; they are not offered as
historical-priority claims.
